Graph-based Artificial Intelligence (GraphAI) encompasses AI problems formulated using graphs, operating on graphs, or relying on graph structures for learning. Contemporary artificial intelligence research is exploring how structured knowledge from graphs can enhance existing approaches to meet the real world's demands for transparency, explainability and performance.

Characterizing GraphAI performance is challenging because different combinations of graph abstractions, representations, algorithms, and hardware acceleration techniques can trigger unpredicTable changes in efficiency. Although benchmarks enable testing different GraphAI implementations, most cannot currently capture the interaction between effectiveness and efficiency, especially across dynamic graphs and structured knowledge graphs.

This work proposes an empirical approach to GraphAI benchmarking, providing a method that enables experimentally trading off between effectiveness and efficiency across different combinations of graph abstractions, representations, algorithms, and hardware accelerators. This approach defines a taxonomy of GraphAI problems, develops a framework for problem-centric benchmarking that does not require invasive workload analyses or code instrumentation, and provides a methodology for selecting benchmark components based on specific use cases. 

Applying this methodology to the problem of Population Segmentation using Community Detection demonstrates how to construct a benchmark that enables experimentally determining the input, implementation, and system configurations that drive the trade-off between execution time and correctness of community composition. The community detection case study illustrates how systematically constructing a GraphAI benchmark exposes the interaction between input and implementation in a GraphAI problem and offers a blueprint for extending this methodology to future GraphAI problems.



----- Nicole Version ----

Graph-based Artificial Intelligence (GraphAI) is the subfield of AI that focuses on the learning problems that operate on graphs or can otherwise involve graph processing. Characterizing GraphAI performance is challenging because different graph abstractions, representations, algorithms, and hardware acceleration techniques can trigger unpredicTable changes in efficiency when combined together. Although benchmarks enable testing different GraphAI implementations, most currently lack the ability to capture the interaction between effectiveness and efficiency, especially across dynamic graphs and structured knowledge graphs.

This work proposes an empirical approach to GraphAI benchmarking, providing a method that enables experimentally trading off between effectiveness and efficiency across different combinations of graph abstractions, representations, algorithms, and hardware accelerators. This approach defines a taxonomy of GraphAI problems, develops a framework for problem-centric benchmarking that does not require invasive workload analyses or code instrumentation, and provides a methodology for selecting benchmark components based on specific use cases. 

Applying this methodology to the problem of Population Segmentation using Community Detection demonstrates how to construct a benchmark that enables experimentally determining the input, implementation, and system configurations that drive the trade-off between execution time and correctness of community composition. The community detection case study illustrates how systematically constructing a GraphAI benchmark exposes the interaction between input and implementation in a GraphAI problem and offers a blueprint for extending this methodology to future GraphAI problems.


----- Original Version -----

Graph-based Artificial Intelligence (GraphAI) encompasses AI problems formulated using graphs, operating on graphs, or relying on graph structures for learning. Contemporary artificial intelligence research is exploring how structured knowledge from graphs can enhance existing approaches to meet the real world's demands for transparency, explainability and performance.

Characterizing the efficiency and effectiveness of GraphAI is critical for real-world deployment and continuing research but challenging because of the breadth of graph abstractions, representations, algorithms and hardware acceleration techniques developed to accommodate increasingly large and complex graphs. These components create variance individually, and combining them yields unpredicTable changes in the demands on computing resources. Existing benchmarks partially address these challenges but fall short of covering the complexity of modern GraphAI workloads and metrics that focus on the effectiveness or efficiency of a workload but rarely cover how they interact.

This work proposes an empirical approach to construct a comprehensive GraphAI benchmark to give evaluators the power to experimentally trade-off between effectiveness and efficiency across many GraphAI problems, with different combinations of graph abstractions, representations, algorithms and hardware accelerators. It defines a taxonomy of GraphAI problems, develops a framework for problem-centric benchmarking and provides a methodology for selecting tasks, datasets, and metrics based on specific use cases. Applying the methodology to the problem of Population Segmentation using Community Detection demonstrates how to create a benchmark that enables an evaluator to quickly try different configurations with experimental results that show the impacts of trading execution time for the correctness of community composition.

The case study illustrates how systematic benchmark construction is beneficial for evaluating different approaches to Community Detection. It offers a blueprint for adding future problems to the proposed GraphAI benchmarking framework. 